% §5 Contact-Centre KPI-Tuned Evaluation Protocol on Public Benchmarks
This section specifies how the trajectory grammar and the stopping rule are evaluated on public task-oriented dialogue benchmarks.
The protocol is tuned to contact-centre KPIs: every terminal class in the outcome set $T = \{\theta_{cont}, \theta_{aban}, \theta_{hand}, \theta_{esc}\}$ of Section~\ref{sec:grammar} --- contained, abandoned, handoff, compliance-escape --- is wired to a metric that a contact-centre operator already tracks.
Figure~\ref{fig:eval_stack} summarises the wiring from grammar states and terminal classes to contact-centre telemetry and per-turn decision signals.
All experimental evidence comes from public datasets; the KPI wiring below is the interface by which public outcome labels acquire operational meaning.

\begin{figure}[t]
  \centering
  \resizebox{\linewidth}{!}{%
    % fig_eval_stack.tex — banking KPI wiring: trajectory classes to telemetry and interventions
\begin{tikzpicture}[
  font=\small,
  box/.style={rectangle, rounded corners, draw, minimum height=0.75cm, inner sep=5pt, align=center},
  kpi/.style={box, fill=olive!8},
  layer/.style={box, fill=blue!6},
  act/.style={box, fill=red!6},
  every edge/.append style={-Stealth, font=\scriptsize},
  node distance=0.55cm and 0.8cm,
  
]
% sensing layer
\node[layer] (sens) {sensing: partial transcript + turn states + tool events + endpointing/barge-in flags};
% grammar layer
\node[layer, below=of sens] (gram) {trajectory grammar $G=(Q,\Sigma,\delta)$: state estimate + cycle depth $D$};
% stopping layer
\node[layer, below=of gram] (stop) {stopping rule over $P(\theta\mid x_{1:k})$: decide at $k_{\mathrm{dec}}\leq k^\ast$, confidence $1-\delta$};
% KPI classes
\node[kpi, below left=0.8cm and -1.2cm of stop] (cont) {contained};
\node[kpi, right=of cont] (aban) {abandoned};
\node[kpi, right=of aban] (hand) {handoff};
\node[kpi, right=of hand] (comp) {compliance-escape};
% interventions
\node[act, below=1.5cm of cont] (none) {continue (no action)};
\node[act, below=1.5cm of aban] (probe) {proactive clarification};
\node[act, below=1.5cm of hand] (esc) {warm transfer to teller};
\node[act, below=1.5cm of comp] (fraud) {fraud flag + policy route};

\draw (sens) -- (gram);
\draw (gram) -- (stop);
\draw (stop.south) -- ++(0,-0.35) -| (cont.north);
\draw (stop.south) -- ++(0,-0.35) -| (aban.north);
\draw (stop.south) -- ++(0,-0.35) -| (hand.north);
\draw (stop.south) -- ++(0,-0.35) -| (comp.north);
\draw (cont) -- (none);
\draw (aban) -- (probe);
\draw (hand) -- (esc);
\draw (comp) -- (fraud);
% latency accounting side note
\node[box, fill=gray!8, right=1.0cm of stop.east, align=left, font=\scriptsize]
  (lat) {decision latency $\leq \Lambda^\ast$ headroom\\ intervention clears cessation floor};
\draw[dashed, gray] (stop) -- (lat);
\end{tikzpicture}%
  }
  \caption{Contact-centre KPI wiring of the trajectory grammar.
           Terminal classes (contained, abandoned, handoff, compliance escape) map one-to-one
           onto operational service metrics; the trajectory monitor consumes partial transcripts,
           tool events, and turn flags; every intervention action must clear the cessation floor.}
  \label{fig:eval_stack}
\end{figure}

\subsection{KPI-to-Terminal-Class Mapping}
\label{subsec:classes}
The terminal classes of the grammar are formalisations of existing operational metrics, not new quantities.
Containment rate counts conversations terminating in $\theta = \text{contained}$, the resolution-without-human-help target tracked in industrial deployments \citep{nubank2026}.
Abandonment rate counts $\theta = \text{abandoned}$: the customer disconnects before resolution.
Abandonment has been modelled as a queueing phenomenon, estimated from waiting data rather than from the conversation itself \citep{silentabandonment2024}.
Handoff rate counts $\theta = \text{handoff}$, the transfer to a human agent; queue-level control optimises this rate from outside the conversation \citep{queuecontrol2026}.
Compliance-escape rate counts $\theta = \text{compliance\_escape}$: unauthorised commitments, disclosure omissions, or policy violations established by post-session audit.
The grammar moves abandonment and escalation inside the conversation: both become predictable at turn granularity, converting reporting KPIs into actionable ones.

\subsection{Public Benchmarks and Label Instantiation}
\label{subsec:datasets}
The protocol instantiates on three public benchmarks, chosen for label structure, turn-level annotation, and comparability with published early-prediction baselines.

\textbf{MultiWOZ~2.4} (primary) provides multi-domain task-oriented dialogues with per-turn belief states.
Following the partial-dialogue forecasting protocol of \citet{nourbakhsh2026tracer}, every visible prefix of a dialogue becomes one supervised instance whose label is the dialogue's terminal outcome, so a dialogue of length $L$ contributes $L$ instances sharing one label.
The terminal instantiation is \emph{binary}: task success versus failure, derived from goal satisfaction in the final dialogue state, exactly as in the TRACER protocol whose public implementation we reuse for comparability \citep{nourbakhsh2026tracer}.
This is the benchmark on which calibration anchors are published: EDS 0.727 at AUC 0.741 under oracle belief states, with the trigger threshold frozen on the development set under a false-positive-rate constraint of $0.15$ \citep{nourbakhsh2026tracer}.

\textbf{ABCD} (customer-service transfer) provides human--human customer-support dialogues (order management, returns, account actions) with an explicit per-turn agent-action taxonomy, the closest public analogue of grammar states.
It is used as a zero-shot transfer target in the spirit of \citet{nourbakhsh2026tracer}, who report the strongest cross-dataset transfer on it (AUC 0.771), and for a leave-one-domain-out validation of the grammar features.

\textbf{$\tau$-bench} (tool-native outcomes) evaluates tool-using customer-service agents in retail and airline domains with environment-verified terminal outcomes: pass/fail by action match and database state, plus a write-check that flags unauthorised environment modifications \citep{sierra2024taubench}.
It is the only public source with native per-turn tool calls and environment-audited outcomes, so the tool-loop channel and the compliance analogue are measured here rather than imputed; trajectories are agent-generated with simulated users, a distributional caveat we state rather than hide.

\emph{Honest scope of the outcome set.}
No public dataset carries all four terminal classes: MultiWOZ and ABCD support the binary instantiation, MultiWOZ's aborted-dialogue flag provides a weak abandoned-like subclass, and $\tau$-bench's write-check is an agent-behaviour compliance analogue rather than a disclosure audit.
The four-class grammar remains the paper's formal object and its operational interface (Subsection~\ref{subsec:classes}); the public evaluation validates the stopping machinery and the drift statistic on the binary instantiation, and the handoff and compliance-escape classes are validated only through these analogues.
State-of-the-art claimed on the four-class problem would require production telemetry; we do not claim it here.

\subsection{Metrics}
Two metrics are reused with credit from TRACER \citep{nourbakhsh2026tracer}: the Early Detection Score, $\mathrm{EDS}_i = (L_i - t_i^*)/L_i$ with $L_i$ the dialogue length in turns, and the mean detection turn $\bar{t}^{\,*}$, computed per terminal class rather than for a single binary outcome.
For calibration, TRACER reports EDS 0.727 at AUC 0.741 under oracle belief states on MultiWOZ \citep{nourbakhsh2026tracer}; these numbers anchor the regime our protocol extends, and our MultiWOZ runs are directly comparable to them by construction.
The cost-weighted detection latency extends both with the operational cost model, weighting each misclassification by how early it was committed:
\begin{equation}
  \mathcal{L}_{\mathrm{cwe}} \;=\; \sum_{\theta \in T \setminus \{\theta_{esc}\}} \pi_\theta \;
  \mathbb{E}_{k_{\mathrm{dec}}}\!\left[\, c(\hat{\theta}_{k_{\mathrm{dec}}}, \theta) \cdot
  w(k_{\mathrm{dec}}) \,\right],
  \qquad w(k) = (H - k)/H,
  \label{eq:cost_weighted_error}
\end{equation}
mirroring the early-prediction reward that TRACER uses at training time \citep{nourbakhsh2026tracer}.
Three definitions make the metric runnable.
The normalisation uses a fixed horizon $H$ (the benchmark's 90th-percentile dialogue length), not the realised length, so the weight is computable at decision time and the length bias shared with EDS is bounded.
The prior $\pi_\theta$ is the per-fold empirical class distribution, so the metric tracks each test fold's operating mix rather than the global one.
And the metric is a \emph{cost-weighted error} measure ($\mathcal{L}_{\mathrm{cwe}}$): correct decisions contribute zero regardless of timing, which is why decision-time quantiles are reported alongside it and baseline dominance is claimed at per-class matched false-alarm budgets, never pooled ones.
The cost model $c(\hat{\theta}, \theta)$ encodes the operational asymmetries with published constants (Table~\ref{tab:costmodel}): a false containment (predicting resolution while the session was abandoned or escaped) costs more than a false handoff, which only spends agent-minutes.
Compliance escape is excluded from the scalar because its cost is unbounded; it is reported separately as detection rate at audit-precision, never averaged in.
Detection-latency distributions are reported as p50/p90/p99, matching the dual-environment latency reporting of the program's prior work.
Calibration of the stopping rule is checked empirically, per class and per split (never pooled, or the common class absorbs the check): the realised error rate at confidence level $1-\varepsilon$ must not exceed $\varepsilon$ with exact binomial bounds, reliability diagrams are reported per class, and the realised $\mathbb{E}[k_{\mathrm{dec}}]$ of the SPRT variant must meet the bound $K(\varepsilon, c)$ of the Trajectory Stopping Theorem.
For the confidence-sequence variant, which carries no closed-form time bound, decision-time distributions are reported as empirical observations rather than checked against $K$.
Every reported metric states which variant (SPRT or confidence sequence) produced it.

\begin{table}[t]
\centering
\caption{Cost-model constants used in $\mathcal{L}_{\mathrm{cwe}}$ and in the stopping
         boundaries (operational accounting values; relative to a false handoff $= 1$).
         The escape row is a floor, not a finite price.}
\label{tab:costmodel}
\small
\begin{tabular}{lll}
\toprule
Error (declared $\to$ truth) & Cost & Provenance \\
\midrule
false handoff (handoff $\to$ contained) & 1 (reference) & human-agent minutes \\
false containment (contained $\to$ abandoned) & 4 & repeat-contact + handling-time accounting \\
false containment (contained $\to$ handoff-truth) & 2 & delayed-resolution accounting \\
false escape flag (escape $\to$ contained truth) & 0.5 & friction/trust-repair accounting \\
missed escape (contained $\to$ escape truth) & $\infty$ (floor 20) & policy-violation event \\
\bottomrule
\end{tabular}
\end{table}

\subsection{Decision-Latency Accounting}
The trajectory monitor is itself a pipeline stage: per turn it consumes the interim transcript, the tool-event log, and the turn flags, and updates the posterior and the sequential test statistic.
If the update runs on the critical path, the published onset floor inflates to $\Lambda^* + \lambda_{\mathrm{mon}}^{\mathrm{irr}}$, where $\lambda_{\mathrm{mon}}^{\mathrm{irr}}$ is the monitor's irreducible per-turn decision latency (its non-hideable serial component, in the reducible/irreducible decomposition of the latency-floor paper), and the competitiveness condition $\Lambda^* + \lambda_{\mathrm{mon}}^{\mathrm{irr}} \leq \tau_{\mathrm{human}}$ must hold, where $\tau_{\mathrm{human}} \approx 200$~ms is the human conversation clock that the onset floor is bound against \citep{garcia2026latencyfloor}.
Placement off the critical path hides the update under the user's speech, provided the SPRT decision is ready before the next intervention point.
The sensing substrate is already measured in the literature: a published interim-transcript pipeline decides in p99 178--268~ms in the cloud configuration \citep{garcia2026asking}, which bounds the per-turn observation update.
Interventions inherit a second floor: taking the floor mid-drift is a barge-in against the agent's own audio, so every intervention must clear the cessation floor (detect, decide, cease, resume) \citep{garcia2026cessationfloor}.
The protocol therefore reports decision latency in both currencies: in turns, the stopping theorem's unit, and in milliseconds, the floors' unit, joined by the per-turn floor arithmetic.
The $k$-step ReAct bound \citep{garcia2026latencyfloor} doubles as a drift alarm threshold in evaluation: any run segment exceeding $k_{\mathrm{cyc}}$ \texttt{tool-call} iterations is scored off-path, and the alarm's own decision latency is measured against the same floors.

\subsection{Baselines}
Five baselines cover the design space without conversation-internal trajectory structure.
The no-intervention and fixed-schedule triggers reproduce TRACER's heuristic baseline families \citep{nourbakhsh2026tracer}.
A TRACER-style dual-stream predictor, retrained for multi-class outcomes under the same prefix protocol from the authors' public implementation \citep{nourbakhsh2026tracer}, is the nearest learned baseline; its hyperparameters, backbone, and feature spec ship in the release configs.
Queue-only escalation decides from queue state alone, without any conversation-internal signal \citep{queuecontrol2026}.
An EarlyEval-style offline predictor, trained on completed episodes and applied as a live trigger, tests whether offline early prediction transfers to in-conversation decisioning \citep{earlyeval2026}.
The stopping rule must dominate all five under the cost-weighted error metric $\mathcal{L}_{\mathrm{cwe}}$ at per-class matched false-alarm budgets, not merely under ranking metrics.

\subsection{Ablations}
The emissions-channel ablation compares text-only emissions against text plus turn states \citep{x2turn2026}, plus tool events, plus endpointing and barge-in flags.
The motivation is measured rather than assumed: text-only turn-end decisions sit near chance on hard cuts (0.40--0.54 accuracy) \citep{garcia2026asking}, so the ablation quantifies which channel carries the trajectory signal.
The drift-statistic ablation replaces cycle depth $D(\sigma)$ with raw conversation length, and must degrade early prediction if the grammar's claim holds: off-path structure, not duration, is the signal.
The stopping-rule ablation compares TRACER's fixed threshold, $\tau = 0.7788$ under an FPR $\leq 0.15$ constraint \citep{nourbakhsh2026tracer}, against the SPRT at matched error, isolating the value of the sequential rule.

\subsection{Reproducibility}
Every dataset, baseline, and metric in this section is public and re-runnable end to end.
The release ships the evaluation code with pinned versions of MultiWOZ~2.4, ABCD, and $\tau$-bench; the prefix-instance builders; the split manifests as SHA-256 checksums; the baseline configs for the TRACER retraining; and the cost-model constants of Table~\ref{tab:costmodel} with the $\pi_\theta$ convention, so the headline metric is reproducible without access to any private system.
The nearest learned baseline is likewise re-runnable from source: TRACER's implementation is public at \url{https://github.com/erfan-nourbakhsh/TRACER} \citep{nourbakhsh2026tracer}.
Because the paper's formal object is the four-class grammar while the public instantiation is binary, the release also includes the label-derivation scripts that map each benchmark's native signals to the binary instantiation, making the honest-scope boundary of Subsection~\ref{subsec:datasets} checkable rather than asserted.
